% Generated by scripts/assemble.py; corrections belong in source/corrections.json.
\examhead{33}{首都师范大学 2026 年数学分析试卷}{36}
% p037-b11; source PDF page 37
\par\Needspace{5\baselineskip}
\blockanchor{p037-b11}
\majorquestion{一}(24 分) 求极限.\par
% p037-b12; source PDF page 37
\blockanchor{p037-b12}
(1) (8 分)  $\lim_{x\to0}\frac{\tan x-x}{\sin^{3}x}.$\par
% p037-b13; source PDF page 37
\blockanchor{p037-b13}
(2) (8 分) $\lim_{n\to \infty}\sum_{k = 1}^{n}\frac{k\sin\left(\frac{k}{n}\right)}{kn + 1}$\par
% p037-b14; source PDF page 37
\blockanchor{p037-b14}
(3) (8 分) $\lim_{n\to \infty}\left(\frac{\sqrt[n]{a} + \sqrt[n]{b}}{2}\right)^n, a, b > 0.$\par
% p037-b15; source PDF page 37
\par\Needspace{5\baselineskip}
\blockanchor{p037-b15}
\majorquestion{二}(10 分) 证明: $f(x) = x \sin x$ 在 $[0, +\infty)$ 上不一致连续.\par
% p037-b16; source PDF page 37
\par\Needspace{5\baselineskip}
\blockanchor{p037-b16}
\majorquestion{三}(15 分) 解答如下问题:\par
% p037-b17; source PDF page 37
\blockanchor{p037-b17}
(1) 证明: 对任意给定的正整数 $n$, 方程 $x^{n+1} + x^n + \cdots + x = 1$ 在 (0,1) 上存在唯一的根, 记为 $x_n$.\par
% p037-b18; source PDF page 37
\blockanchor{p037-b18}
(2) 证明: 极限 $\lim_{n\to \infty}x_n = \frac{1}{2}$.\par
% p037-b19; source PDF page 37
\par\Needspace{5\baselineskip}
\par\Needspace{6\baselineskip}
\blockanchor{p037-b19}
\majorquestion{四}(12 分) 设\par
\nopagebreak[4]
% p037-b20; source PDF page 37
\blockanchor{p037-b20}
\displayformula{f (x, y) = \left\{ \begin{array}{l l} x - y + \frac {x y ^ {2}}{x ^ {2} + y ^ {2}}, & (x, y) \neq (0, 0); \\ 0, & (x, y) = (0, 0). \end{array} \right.}
% p037-b21; source PDF page 37
\blockanchor{p037-b21}
证明： $f(x,y)$ 在原点 $(0,0)$ 处连续，沿任何方向的方向导数存在，但不可微.\par
% p037-b22; source PDF page 37
\par\Needspace{5\baselineskip}
\par\Needspace{6\baselineskip}
\blockanchor{p037-b22}
\majorquestion{五}(15 分) 设 $f(x)$ 在 $[0, 1]$ 上连续可微, 满足\par
\nopagebreak[4]
% p037-b23; source PDF page 37
\blockanchor{p037-b23}
\displayformula{f ^ {\prime} (x) + f (x) \sin x = \int_ {0} ^ {x} f (t) \mathrm{d} t, x \in [ 0, 1 ].}
% p037-b24; source PDF page 37
\blockanchor{p037-b24}
且 $\int_0^1 f(x)\mathrm{d}x = 0$ ，证明： $f(x)\equiv 0,x\in [0,1]$\par
% p038-b00; source PDF page 38
\par\Needspace{5\baselineskip}
\blockanchor{p038-b00}
\majorquestion{六}(10 分) 证明: 当 $\lambda < 1$ 时, $\lim_{R \to +\infty} R^{\lambda} \int_{0}^{\frac{\pi}{2}} e^{-R \sin \theta} \mathrm{d}\theta = 0$.\par
% p038-b01; source PDF page 38
\par\Needspace{5\baselineskip}
\par\Needspace{6\baselineskip}
\blockanchor{p038-b01}
\majorquestion{七}(12 分) 设 $f(x)$ 在 $[- \pi, \pi]$ 上连续, 且导函数 $f'(x)$ 可积和绝对可积, 若有\par
\nopagebreak[4]
% p038-b02; source PDF page 38
\blockanchor{p038-b02}
\displayformula{f (x) \sim \frac {a _ {0}}{2} + \sum_ {n = 1} ^ {\infty} (a _ {n} \cos n x + b _ {n} \sin n x), x \in [ - \pi , \pi ].}
% p038-b03; source PDF page 38
\par\Needspace{6\baselineskip}
\blockanchor{p038-b03}
证明：\par
\nopagebreak[4]
% p038-b04; source PDF page 38
\blockanchor{p038-b04}
\displayformula{f ^ {\prime} (x) \sim \frac {c}{2} + \sum_ {n = 1} ^ {\infty} [ (n b _ {n} + (- 1) ^ {n} c) \cos n x - n a _ {n} \sin n x ], \dot {x} \in [ - \pi , \pi ].}
% p038-b05; source PDF page 38
\blockanchor{p038-b05}
其中 $c = \frac{f(\pi) - f(-\pi)}{\pi}$, 且 $c = \lim_{n\to \infty}[-1)^{n - 1}nb_n]$.\par
% p038-b06; source PDF page 38
\par\Needspace{5\baselineskip}
\par\Needspace{6\baselineskip}
\blockanchor{p038-b06}
\majorquestion{八}(12 分) 设 $C$ 是单位圆周 $x^{2} + y^{2} = 1$, 方向为逆时针, 求积分\par
\nopagebreak[4]
% p038-b07; source PDF page 38
\blockanchor{p038-b07}
\displayformula{\oint_ {C} \frac {(x - y) \mathrm{d} x + (x + 4 y) \mathrm{d} y}{x ^ {2} + 4 y ^ {2}}.}
% p038-b08; source PDF page 38
\par\Needspace{5\baselineskip}
\par\Needspace{6\baselineskip}
\blockanchor{p038-b08}
\majorquestion{九}(10 分) 计算曲面积分\par
\nopagebreak[4]
% p038-b09; source PDF page 38
\blockanchor{p038-b09}
\displayformula{\iint_ {S} x ^ {2} \mathrm{d} y \mathrm{d} z + y ^ {2} \mathrm{d} z \mathrm{d} x + z ^ {2} \mathrm{d} x \mathrm{d} y.}
% p038-b10; source PDF page 38
\blockanchor{p038-b10}
其中 $S$ 为球面 $(x - a)^2 + (y - b)^2 + (z - c)^2 = R^2 (R > 0)$ 的外侧.\par
% p038-b11; source PDF page 38
\par\Needspace{5\baselineskip}
\par\Needspace{6\baselineskip}
\blockanchor{p038-b11}
\majorquestion{十}(10 分) 设函数列 $f_n(x) \geq 0 (n = 1,2,\dots)$ 与函数 $g(x)$ 在区间 $[a,b]$ 上有界且黎曼可积, 对任意的 $c \in (a,b)$, 当 $n \to \infty$ 时, $\{f_n(x)\}$ 在区间 $[c,b]$ 上一致收敛于 0 , 且\par
\nopagebreak[4]
% p038-b12; source PDF page 38
\blockanchor{p038-b12}
\displayformula{\lim _ {n \rightarrow \infty} \int_ {a} ^ {b} f _ {n} (x) \mathrm{d} x = 1, \lim _ {x \rightarrow a ^ {+}} g (x) = A.}
% p038-b13; source PDF page 38
\blockanchor{p038-b13}
证明： $\lim_{n\to \infty}\int_a^b g(x)f_n(x)\mathrm{d}x = A.$\par
% p038-b14; source PDF page 38
\par\Needspace{5\baselineskip}
\blockanchor{p038-b14}
\majorquestion{十一}(10 分) 设积分 $\int_{0}^{+\infty} f(x) \, \mathrm{d}x$ 收敛, 且存在 $A > 0$, 使得 $f(x)$ 在 $[A, +\infty)$ 上非负, 存在 $0 < p < 1$ 使得当 $x \to +\infty$ 时, $x^p f(x)$ 单调递减趋于零, 证明: $\lim_{x \to +\infty} x f(x) = 0$.\par
% p038-b15; source PDF page 38
\par\Needspace{5\baselineskip}
\blockanchor{p038-b15}
\majorquestion{十二}(10 分) 设  $f_{n}(x)\}$  ( $n=1,2,\cdots$ ) 为 [a,b] 上的连续函数列, 且对任意给定的  $x\in[a,b]$ ,  $\{f_{n}(x)\}$  关于 n 单调递增. 证明: 如果函数列  $\{f_{n}(x)\}$  在 [a,b] 上逐点收敛于连续函数  $f(x)$ , 则  $\{f_{n}(x)\}$  在 [a,b] 上必一致收敛于  $f(x)$ .\par
